%!TEX TS-program = xelatex
%!TEX encoding = UTF-8 Unicode

% 中文译注本主文件：tarasov-qasp-zh.tex
% 英文原工程一个不动；中文各章放 zh/ 目录，用 \InputIfFileExists 按需载入。
% 编译：xelatex（两遍，出目录）

\documentclass[b5paper]{tufte-book}

\usepackage{ifxetex}

% ---- 中文支持（xelatex + ctex/xeCJK，不改 tufte 标题体例） ----
% 注：英文原工程的 mathdesign/tgadventon 是 pdflatex 字体方案，
% 在 xelatex(TU 编码)下无定义；中文版改用 fontspec 加载 TeX Gyre OpenType。
\usepackage[UTF8,heading=false,fontset=windows]{ctex}
\setmainfont{TeX Gyre Pagella}
\setsansfont{TeX Gyre Adventor}

% Book metadata

\title[中学物理问答]{中学物理问答}
\author[塔拉索夫]{列夫·塔拉索夫 \quad 阿尔季娜·塔拉索娃}
\publisher{Mir 出版社，莫斯科}
\date{1973 年初版 \quad·\quad 中文译注本}

% For nicely typeset tabular material
\usepackage{booktabs}

% For graphics / images
\usepackage{graphicx}
\setkeys{Gin}{width=\linewidth,totalheight=\textheight,keepaspectratio}

\usepackage{fancyvrb}
\fvset{fontsize=\normalsize}

\newcommand{\hangp}[1]{\makebox[0pt][r]{(}#1\makebox[0pt][l]{)}}
\newcommand{\hangstar}{\makebox[0pt][l]{*}}
\usepackage{xspace}

\newcommand{\openepigraph}[2]{%
  \begin{fullwidth}
  \sffamily\large
  \begin{doublespace}
  \noindent\allcaps{#1}\\% epigraph
  \noindent\allcaps{#2}% author
  \end{doublespace}
  \end{fullwidth}
}

\newcommand{\blankpage}{\newpage\hbox{}\thispagestyle{empty}\newpage}

\newcommand{\hlred}[1]{\textcolor{Maroon}{#1}}% prints in red
\newcommand{\hlblue}[1]{\textcolor{SteelBlue}{#1}}% prints in blue
\newcommand{\drkgry}[1]{\textcolor{darkgray}{{#1}}}
\newcommand{\hangleft}[1]{\makebox[0pt][r]{#1}}
\newcommand{\hairsp}{\hspace{1pt}}
\newcommand{\hquad}{\hskip0.5em\relax}
\newcommand{\TODO}{\textcolor{red}{\bf TODO!}\xspace}
\newcommand{\ie}{\textit{i.\hairsp{}e.}\xspace}
\newcommand{\eg}{\textit{e.\hairsp{}g.}\xspace}
\newcommand{\na}{\quad--}
\newcommand{\tuftebs}{\symbol{'134}}

\usepackage{pdfpages}

% titles formatting

\titleformat{\chapter}
  {\normalfont\itshape\Huge\color{darkgray}}
  {\thechapter}{1em}{}
\titleformat{\section}
  {\normalfont\LARGE\scshape\color{IndianRed}}
  {\thesection}{1em}{}

% math related pacakges

\usepackage{amsmath}
\usepackage{mathabx}

% for nice lists including inline lists
\usepackage[inline]{enumitem}

\setlength\parskip{1em plus 0.1em minus 0.2em}

\newlength{\oldparindent}
\setlength{\oldparindent}{\parindent}
\newcommand{\myindent}{\hspace{\oldparindent}}

\setlength{\parindent}{0pt}

% for setting the dialogues
\usepackage{dialogue}

% 中文对话角色标签：粗体 + 全角冒号（smallcaps 不适用于中文）
\renewcommand*\DialogueLabel [1] {\textbf{#1}：}

% custom commands for speaking

\newcommand{\tchr}{\speak{教师}}
\newcommand{\std}{\speak{学生}}
\newcommand{\stda}{\speak{学生甲}}
\newcommand{\stdb}{\speak{学生乙}}
\newcommand{\stdc}{\speak{学生丙}}

% command for refering to a figure easily use \figr{ref}
\newcommand{\figr}[1] {\textcolor{darkgray}{图~\ref{#1}}}

% command for refering to a chapter
\newcommand{\chap}[1] {\textcolor{darkgray}{第~\ref{#1}~章}}

% 译注层评注宏：边栏【评注】。这是本译注本的核心体例，冻结于样板章定稿。
\newcommand{\pnote}[1]{\marginnote{\textbf{【评注】}\;#1}}

% table of contents formatting

\usepackage{tocloft}
\renewcommand{\cfttoctitlefont}{\hfill\Huge\itshape\flushleft\color{darkgray}}
\renewcommand{\cftchapfont}{\Large\itshape\color{darkgray}}
\renewcommand\cftsecfont{\color{IndianRed}\large\itshape}

\renewcommand{\cftchappagefont}{\Large\color{SteelBlue}}
\renewcommand{\cftsecpagefont}{\large\color{IndianRed}}
\renewcommand{\cftchappresnum}{\textsection\,}
\cftsetindents{chapter}{0em}{2.5em}
\cftsetindents{section}{2.5em}{2.5em}
\renewcommand{\cftsecdotsep}{\cftnodots}

\makeatletter
\renewcommand{\@tocrmarg}{3em plus1fil}
\makeatother

% for no headers and footers on empty pages
\usepackage{emptypage}

% headers and footers

\makeatletter
\renewcommand\mainmatter{%
  \cleardoublepage%
  \@mainmattertrue%
  \fancyhf{}%
  \ifthenelse{\boolean{@tufte@twoside}}%
    {% two-side
      \renewcommand{\chaptermark}[1]{\markboth{##1}{}}%
      \fancyhead[LE]{\small\thepage\quad{\newlinetospace{\plaintitle}}}%
      \fancyhead[RO]{\small{\newlinetospace{\leftmark}}\quad\thepage}%
    }%
    {% one-side
      \fancyhead[RE,RO]{\small{\newlinetospace{\plaintitle}}\quad\thepage}%
    }%
}
\makeatother

\setcounter{secnumdepth}{1}

\titleformat{\chapter}
  [block]% shape
  {\relax\ifthenelse{\NOT\boolean{@tufte@symmetric}}{\begin{fullwidth}}{}}% format applied to label+text
  {\color{darkgray}\itshape\huge~\textsection~\thechapter}% label
  {1em}% horizontal separation between label and title body
  {\color{darkgray}\huge\rmfamily\itshape}% before the title body
  [\ifthenelse{\NOT\boolean{@tufte@symmetric}}{\end{fullwidth}}{}]% after the title body

% nice units
\usepackage[per-mode=symbol]{siunitx}
\DeclareSIUnit\calorie{cal}
% for drawing using TikZ

\usepackage{tikz}
\usetikzlibrary{arrows,decorations}
\usetikzlibrary{arrows.meta}
\usetikzlibrary{calc}
\usetikzlibrary{shapes.geometric}
\usetikzlibrary{patterns,calc}
\usetikzlibrary{angles,quotes}
\usetikzlibrary{patterns}
\usetikzlibrary{positioning}
\usetikzlibrary{decorations.markings}
\usetikzlibrary{backgrounds}

\newdimen\XCoord
\newdimen\YCoord
\newcommand*{\ExtractCoordinate}[1]{\path (#1); \pgfgetlastxy{\XCoord}{\YCoord};}%

\tikzset{middlearrow/.style={
        decoration={markings,
            mark= at position 0.5 with {\arrow{#1}} ,
        },
        postaction={decorate}
    }
}

\usepackage{tkz-euclide}

\usepackage{calrsfs}
\usepackage{nicefrac}

\everymath{\displaystyle}

\DeclareMathAlphabet{\pazocal}{OMS}{zplm}{m}{n}
\newcommand{\Ea}{\mathcal{E}}

% for multi columns in list
\usepackage{multicol}

% continuous numbering of equations and figures

\usepackage{chngcntr}
\counterwithout{figure}{chapter}
\counterwithout{equation}{chapter}

\hypersetup{colorlinks,linkcolor=darkgray}

% ---- 中文本地化 ----
\renewcommand{\contentsname}{目录}
\renewcommand{\figurename}{图}
\renewcommand{\tablename}{表}
% CJK 不适用 smallcaps/allcaps（大写与 letterspace 在 xelatex 下不可用），改用粗体
\renewcommand{\smallcaps}[1]{\textbf{#1}}
\renewcommand{\allcaps}[1]{\textbf{#1}}

% [zh 版已知问题] 当前 MiKTeX 25.4 的 LaTeX 格式与 ltcmd 钩子系统存在版本偏差：
% tufte 浮动体在 \end{figure} 构建图题盒时，hyperref 的浮动锚点路径会触发
% 内核包装版 \label 中未定义的 \__hook_toplevel 引用，每个图产生 3 个可恢复错误
% （Undefined control sequence + 两条 minipage 连锁）。输出与全部交叉引用实际正确；
% 英文原工程在同一环境下重编也会出现同样错误（上游问题，非译本内容问题）。
% 已验证无效的尝试：AddToHook{label}{} 注册、自定义 \__hook_toplevel 分发器、
% tufte-common.def 安全发射器（后者保留，无副作用）。若日后 MiKTeX 修复格式，
% 此处错误应自动消失。

\begin{document}

% !TEX root = tarasov-qasp-zh.tex
\thispagestyle{empty}

\begin{tikzpicture}[remember picture,overlay]
  \fill[thick,yellow!10]
        (current page.north west) rectangle (current page.south east);
\end{tikzpicture}

\begin{tikzpicture}[remember picture,overlay]
\fill [yellow](5.35,-16.95) circle [radius=.92cm];
\end{tikzpicture}
\begin{tikzpicture}[remember picture,overlay]
\fill [yellow](6.8,-16.95) circle [radius=.92cm];
\end{tikzpicture}
\begin{tikzpicture}[remember picture,overlay]
\fill [yellow](8.2,-16.95) circle [radius=.92cm];
\end{tikzpicture}
\begin{tikzpicture}[remember picture,overlay]
\fill [red](10.8,-15.45) circle [radius=.45cm];
\end{tikzpicture}
\begin{tikzpicture}[remember picture,overlay]
\fill [red](2.43,-15.45) circle [radius=.45cm];
\end{tikzpicture}
\begin{tikzpicture}[remember picture,overlay]
\fill [red](10.5,-5.65) circle [radius=.45cm];
\end{tikzpicture}
\begin{tikzpicture}[remember picture,overlay]
\fill [red](2.3,-5.68) circle [radius=.45cm];
\end{tikzpicture}

\begin{tikzpicture}[remember picture,overlay]
\draw [in=180,out=270,looseness=.8,ForestGreen, line width=6pt](2.78,-5.3) to (10.6,-14.65);
\end{tikzpicture}
\begin{tikzpicture}[remember picture,overlay]
\draw [in=270,out=0,looseness=.8,ForestGreen, line width=6pt](3,-14.65) to (10.8,-5.3);
\end{tikzpicture}

\begin{tikzpicture}[remember picture,overlay]
\node [scale=3.,yshift=-1.75cm,xshift=-0.85cm]	 at (current page.center) {$m$};
\end{tikzpicture}

\begin{tikzpicture}[remember picture,overlay]
\node [scale=3.,yshift=-1.75cm,xshift=0.85cm]	 at (current page.center) {$a$};
\end{tikzpicture}

\begin{tikzpicture}[remember picture,overlay]
  \tkzDefPoints{6.95/-5.4/A}
  \tkzDrawArc[R,line width=0.5pt,color=red](A,3.95 cm)(35,145);
  \tkzDrawArc[R,line width=0.6pt,color=red](A,3.48 cm)(32.75,147.5);
  \tkzDrawArc[R,line width=0.7pt,color=red](A,3. cm)(29.75,151);
    \tkzDrawArc[R,line width=0.8pt,color=red](A,2.7 cm)(27.5,154);
        \tkzDrawArc[R,line width=0.9pt,color=red](A,2.4 cm)(24.5,157.5);
    \tkzDrawArc[R,line width=1pt,color=red](A,2.2 cm)(22,160);
        \tkzDrawArc[R,line width=1.05pt,color=red](A,2.05 cm)(20,163);
    \tkzDrawArc[R,line width=1.1pt,color=red](A,1.93 cm)(18,166);
    \tkzDrawArc[R,line width=1.15pt,color=red](A,1.83 cm)(16,169);
\end{tikzpicture}

\begin{tikzpicture}[remember picture,overlay]
\node [scale=2.6,yshift=3.6cm,font=\sffamily,darkgray]	 at (current page.center) {列夫·塔拉索夫 \quad 阿尔季娜·塔拉索娃};
\end{tikzpicture}

\begin{tikzpicture}[remember picture,overlay]
\node [scale=3.65,align=center,font=\sffamily,yshift=1.9cm,xshift=0cm,darkgray] at (current page.center) {\textbf{中学物理} \\ \textbf{问\ \ 答}};
\end{tikzpicture}

\begin{tikzpicture}[remember picture,overlay]
\node [yshift=-3.1cm]	 at (current page.center) {\includegraphics[width=0.9\textwidth]{figs/front-back/tarasov-qasp-fc3.pdf}};
\node [scale=2,yshift=-5.5cm,font=\sffamily,darkgray]	 at (current page.center) {MIR 出版社 \quad 莫斯科};
\end{tikzpicture}

\cleardoublepage


\frontmatter

\maketitle

% !TEX root = tarasov-qasp-zh.tex
\thispagestyle{empty}

{\noindent
本书俄文原著由 \emph{Nicholas Weinstein} 译为英文，1973 年出版（据 1968 年俄文修订版）。\\[4pt]
电子重排本由 \emph{Damitr Mazanav} 以 \LaTeX{} 排版，
2020 年经 \url{http://mirtitles.org} 发布于网络。\\
重排本源文件仓库：\url{https://gitlab.com/mirtitles/tarasov-qasp/}\\[10pt]
中文译注本据上述 \LaTeX{} 重排本转译。\\
本译注本仅限个人学习使用，不作公开分发。}

\cleardoublepage


\setcounter{tocdepth}{3}

\begingroup
 \hypersetup{hidelinks}
 \tableofcontents
\endgroup

% !TEX root = tarasov-qasp-zh.tex

\vspace*{1cm}

\chapter[译注本说明]{译注本说明}
\label{zh-note}

本册是塔拉索夫（Lev Tarasov）与塔拉索娃（Aldina Tarasova）《中学物理问答》（\emph{Questions and Answers in School Physics}）的中文译注本。体例如下：

\begin{itemize}
\item \textbf{正文}为原书的完整忠实翻译。教师与学生的对话（包括学生答错、被指出、再纠正的全部过程）原样保留——这种往复正是原书的教学法所在。
\item \textbf{边栏中以【评注】开头的内容为译注层，不是原书文字。}评注只做三类事：对原书为中学水平而从简的表述做精确化与勘误；对原书成书年代（1968/1973）距今的必要说明；与中国中学物理教材（人教版）体系的对接与延伸。评注由 AI 辅助生成，供批判性参考。
\item \textbf{插图}沿用原书扫描件：图题已译为中文，图内文字保留英文原文。
\item \textbf{底本}为 mirtitles.org 2020 年 \LaTeX{} 重排本（社区录入，非出版级校样）。遇可疑录入之处，正文按可核实的正确内容译出，并在评注中说明；无法核实的标注"待考"。
\end{itemize}

\cleardoublepage


% !TEX root = tarasov-qasp-zh.tex

\vspace*{1cm}

\chapter[俄文版编者的话]{俄文版编者的话}
\label{preface}

可以断言，还没有一个正在准备高校物理入学考试的学生，翻开过与本书类似的书。作者采用极其生动的对话形式，全面讨论了教学大纲中几乎所有的课题，特别是那些通常被认为难以理解的问题。书中对考生在物理入学考试中常犯的错误作了详细分析。学生会发现这是一本异常清晰而有趣的读物：它从多个角度处理复杂问题，并配有大量有助于更深入理解相关概念与思想的精彩插图。作者是莫斯科电子技术学院（Moscow Institute of Electronics Engineering）的教师，非常了解报考高校学生的总体训练水平，并有多年主持入学考试的经验。作者的专业学识，加上生动明晰的讲解，使本书成为准备物理考试的学生极为有用的辅导读物。\\
\begin{flushright}
\emph{G. 叶皮法诺夫}（G.~Epifanov）教授，物理数学博士
\end{flushright}

\chapter{作者前言}
\label{foreword}

本书是为准备高校物理入学考试的学生编写的辅助读物，采用作者（\textbf{教师：}）与好问的读者（\textbf{学生：}）对话的形式。这对于分析考生在入学考试中的常见错误、复习同一问题的多种解法、讨论物理理论中的疑难问题，都特别方便。

书中讨论了大量中学物理的问题与习题，此外还配备了供课后自测的习题（附答案）。大部分问题与习题出自莫斯科电子技术学院 1964—1966 年的入学考试。对考生错误的分析总是有教益的：它可以引出问题的不同侧面，点出某些细微之处，从而达到对基础知识更透彻的理解。

\pnote{原书成书于 1968/1973 年，"入学考试"指苏联工科高校的招生考试（口试占重要地位）。对今天的中国读者，本书的合适位置是：学完高中物理后的复习查漏与竞赛入门思维训练；它不是入门教材，作者在下文也是这样定位的。}

不过，这种分析做起来非常困难。正确的答案只有一个，错误的答案却可以有很多。任何问题都不可能预见全部错误答案——许多错误永远藏在口试学生令人着急的沉默背后。尽管如此，某些问题特有的错误答案总是反复出现；有许多问题几乎不可避免地会被答错。本书主要就以这类问题和习题为基础。

我们要提醒读者：本书绝不是一本涵盖大纲全部条目的教材。读者在书中找不到物理课程所要求的系统叙述。本书更像一篇随意讲出的故事，或者不如说，一场随意展开的讨论。因此，对那些想从头学习物理或系统整理物理知识的人，本书用处不大。它是为那些在考试前夕希望扩充物理知识的人准备的。

在我们心目中，理想的读者已经学完中学物理课程，对物理内容有清晰的整体图景，记得主要关系式，能引用各种定律，对所用的单位也相当熟悉。他正处于那种"悬置"状态：不再是中学生，也还不是正式的大学生——但他渴望成为大学生。如果为此需要扩展物理知识，本书可以帮他一把。

首先，我们希望本书能证明：死记硬背教材（哪怕是很好的教材）不仅令人厌倦，而且毫无成效。学生必须学会思考，学会琢磨材料本身，而不是单纯背诵。如果读者能在某种程度上获得这种认识，我们就认为自己的努力是值得的。

最后，我们感谢 G. 叶皮法诺夫教授，没有他的鼓励和宝贵帮助，本书不可能写成并出版。我们还衷心感谢莫斯科动力学院（Moscow Power Engineering Institute）物理系高级教师 E. N. 弗托罗夫（E.~N.~Vtorov）、副教授 A. G. 切尔托夫（A.~G.~Chertov）和 V. A. 法布里坎特教授（V.~A.~Fabrikant）对手稿提出的许多有益建议和建设性批评。

\begin{flushright}
\textit{L. 塔拉索夫}\\
\textit{A. 塔拉索娃}
\end{flushright}

\cleardoublepage


\mainmatter

\InputIfFileExists{zh/ch-01-zh}{}{}
\InputIfFileExists{zh/ch-02-zh}{}{}
\InputIfFileExists{zh/ch-03-zh}{}{}
\InputIfFileExists{zh/ch-04-zh}{}{}
\InputIfFileExists{zh/ch-05-zh}{}{}
\InputIfFileExists{zh/ch-06-zh}{}{}
\InputIfFileExists{zh/ch-07-zh}{}{}
\InputIfFileExists{zh/ch-08-zh}{}{}
\InputIfFileExists{zh/ch-09-zh}{}{}
\InputIfFileExists{zh/ch-10-zh}{}{}
\InputIfFileExists{zh/ch-11-zh}{}{}
\InputIfFileExists{zh/ch-12-zh}{}{}
\InputIfFileExists{zh/ch-13-zh}{}{}
\InputIfFileExists{zh/ch-14-zh}{}{}
\InputIfFileExists{zh/ch-15-zh}{}{}
\InputIfFileExists{zh/ch-16-zh}{}{}
\InputIfFileExists{zh/ch-17-zh}{}{}
\InputIfFileExists{zh/ch-18-zh}{}{}
\InputIfFileExists{zh/ch-19-zh}{}{}
\InputIfFileExists{zh/ch-20-zh}{}{}
\InputIfFileExists{zh/ch-21-zh}{}{}
\InputIfFileExists{zh/ch-22-zh}{}{}
\InputIfFileExists{zh/ch-23-zh}{}{}
\InputIfFileExists{zh/ch-24-zh}{}{}
\InputIfFileExists{zh/ch-25-zh}{}{}
\InputIfFileExists{zh/ch-26-zh}{}{}
\InputIfFileExists{zh/ch-27-zh}{}{}
\InputIfFileExists{zh/ch-28-zh}{}{}
\InputIfFileExists{zh/ch-29-zh}{}{}
\InputIfFileExists{zh/ch-30-zh}{}{}
\InputIfFileExists{zh/ch-31-zh}{}{}
\InputIfFileExists{zh/ch-32-zh}{}{}
\InputIfFileExists{zh/ch-33-zh}{}{}

\backmatter

\InputIfFileExists{zh/ch-34-ans-zh}{}{}
\InputIfFileExists{zh/00-back-zh}{}{}

\end{document}
